\uuid{M5mC}
\titre{Ensembles}
\auteur{bodin ; croizat ; sacre}
\organisation{exo7}
\licence{CC BY-NC-SA 4.0}
\datecreate{2022-03-08}
\niveau{L1}
\module{Algèbre}
\chapitre{Logique, ensemble, raisonnement}
\sousChapitre{Ensemble}
\difficulte{3}
\theme{Ensembles}
\exo7qcmid{hilisit.logique.036}

\contenu{
\begin{question}
\qclassification{100.02}
Soient $A$ et $B$ deux parties d'un ensemble $E$.
Quelles sont les affirmations vraies (quel que soit le choix de $A$ et $B$) ?
\begin{answers}
	\good{$(A\cap B)\cup(A\setminus B)=A$}
	\bad{$(A\cap B)\cup(A\setminus B)=B$}
	\bad{$(A\cap B)\cup(A\setminus B)= A\cup B$}
	\good{$(A\cap B)\cup(A\setminus B)\cup(B\setminus A)=A\cup B$}
\end{answers}
\begin{explanations}
Si l'on réunit les éléments qui sont dans $A$ et dans $B$ avec ceux qui sont dans $A$ et pas dans $B$, on obtient tous les éléments de $A$.\\
On a 3 possibilités pour les éléments qui sont dans $A$ ou dans $B$ : ils sont dans $A$ et pas dans $B$, ou dans $A$ et dans $B$, ou dans $A$ et dans $B$.
Faites des schémas avec des "patates" (diagrammes de Venn) pour visualiser ces propriétés.
\end{explanations}
\end{question}
}
